\lysh{14961}{4}

\begin{alist}
\item Η γραφική παράσταση της συνάρτησης $f(x) = x^2 - x - 1$ τέμνει τον άξονα $x'x$
στα σημεία $A(\omega, 0)$, $B(\phi, 0)$, οπότε $\omega, \phi$ είναι οι ρίζες της εξίσωσης $f(x) = 0$ με
$\omega < 0 < \phi$ αφού το σημείο $O(0, 0)$ είναι μεταξύ των $A$ και $B$ στον άξονα $x'x$.

Η εξίσωση $f(x) = 0 \Leftrightarrow x^2 - x - 1 = 0$ έχει διακρίνουσα $\Delta = 5$ και ρίζες τους
αριθμούς $\dfrac{1 + \sqrt{5}}{2}$ και $\dfrac{1 - \sqrt{5}}{2}$. Επειδή $\dfrac{1 - \sqrt{5}}{2} < 0 < \dfrac{1 + \sqrt{5}}{2}$ έχουμε ότι $\phi = \dfrac{1 + \sqrt{5}}{2}$ και $\omega = \dfrac{1 - \sqrt{5}}{2}$.

\begin{rlist}
\item Αφού $\omega, \phi$ είναι οι ρίζες της εξίσωσης $f(x) = 0$ δηλαδή της $x^2 - x - 1 = 0$
έχουμε ότι $\omega + \phi = -\dfrac{\beta}{\alpha} = -\dfrac{-1}{1} = 1$.

\item Αφού $\omega, \phi$ είναι οι ρίζες της εξίσωσης $f(x) = 0$ δηλαδή της $x^2 - x - 1 = 0$
έχουμε ότι $\omega \cdot \phi = \dfrac{\gamma}{\alpha} = \dfrac{-1}{1} = -1$.
\end{rlist}

\item Είναι $(OB) = |\phi| = \left|\dfrac{1 + \sqrt{5}}{2}\right| = \dfrac{1 + \sqrt{5}}{2}$ και $(OA) = |\omega| = \left|\dfrac{1 - \sqrt{5}}{2}\right| = \dfrac{\sqrt{5} - 1}{2}$.
Επειδή $\dfrac{\sqrt{5} + 1}{2} > \dfrac{\sqrt{5} - 1}{2}$ συμπεραίνουμε ότι $(OB) > (OA)$.

Εναλλακτικά, είναι $(OB) = |\phi| = \phi$ αφού $\phi > 0$ και $(OA) = |\omega| = -\omega$ αφού $\omega < 0$.
Είναι $(OB) > (OA) \Leftrightarrow \phi > -\omega \Leftrightarrow \phi + \omega > 0 \Leftrightarrow 1 > 0$ που ισχύει.

\item Αφού ο $\beta$ είναι μεγαλύτερος από τον αντίστροφό του και η διαφορά τους
ξεπερνάει τη μία μονάδα, έχουμε ότι $\beta - \dfrac{1}{\beta} > 1$ και εφόσον $\beta > 0$ έχουμε ισοδύναμα
$\beta^2 - 1 > \beta \Leftrightarrow \beta^2 - \beta - 1 > 0$. Το τριώνυμο $x^2 - x - 1$ έχει ρίζες $\omega, \phi$ και
γίνεται θετικό, δηλαδή ομόσημο του $\alpha = 1$, για $x < \omega$ ή $x > \phi$. Συνεπώς $\beta < \omega$ ή $\beta > \phi$. Όμως $\beta > 0$, οπότε $\beta > \phi$.

Εναλλακτικά, $\beta^2 - \beta - 1 > 0 \Leftrightarrow f(\beta) > 0$. Από τη γραφική παράσταση της $f$
βλέπουμε ότι τα σημεία της γραφικής παράστασης που έχουν θετική τεταγμένη,
δηλαδή είναι πάνω από τον άξονα $x'x$, είναι αυτά που είναι δεξιά του $B$ ή αριστερά
του $A$. Συνεπώς $\beta < \omega$ ή $\beta > \phi$. Όμως $\beta > 0$, οπότε $\beta > \phi$.

\item Είναι $f\left(\dfrac{5}{3}\right) = \left(\dfrac{5}{3}\right)^2 - \dfrac{5}{3} - 1 = \dfrac{25}{9} - \dfrac{15}{9} - \dfrac{9}{9} = \dfrac{1}{9} > 0$ και επειδή $\dfrac{5}{3} > 0$ έχουμε ότι $\phi < \dfrac{5}{3}$.

Εναλλακτικά, $\dfrac{5}{3} - \dfrac{3}{5} = \dfrac{25}{15} - \dfrac{9}{15} = \dfrac{16}{15} > 1$, οπότε με βάση το $\gamma)$ έχουμε ότι $\phi < \dfrac{5}{3}$.
\end{alist}

